\documentclass[12pt, oneside, final]{report}
\usepackage{geometry}
\geometry{a4paper, left=20mm, right=20mm, top=25mm, bottom=25mm}
\usepackage[utf8]{inputenc}
\usepackage{t1enc}
\usepackage[MeX]{polski}
\usepackage{graphicx}
\usepackage{amsmath}
\usepackage{amssymb}
\usepackage{mathtools}
\usepackage{indentfirst}
\usepackage{pdfpages}
\usepackage{xcolor}
\usepackage{placeins} % provides \FloatBarrier
\usepackage{tikz}
\usetikzlibrary{positioning,shapes,arrows,calc,decorations.markings,shadows}

% Czcionka
\usepackage{charter}

% Pojedyncze elementy na górze stron
\makeatletter
\setlength{\@fptop}{0pt}
\makeatother

% Styl tytułowania rozdziałów:
\usepackage{titlesec}
%\titleformat{\chapter}{\normalfont\huge}{\bf\thechapter.}{20pt}{\huge\bf}
% Styl tytułowania rozdziałów:
\titleformat{\chapter}[display] 
{\centering\normalfont\huge\bfseries}{\centering\chaptertitlename\ \thechapter.}{0.5em}{}
\titlespacing{\chapter}{0em}{0em}{2em}

% dodatkowe kropki po numerach rozdziałów, sekcji itd
\usepackage{titlesec}
\titlelabel{\thetitle.\quad}

% Dodatkowe odstępy w tabelach (do UseCaseow)
\usepackage{array}
\setlength\extrarowheight{4pt}


% Szablon do Use Caseów
\newcounter{UseCaseCounter}
\setcounter{UseCaseCounter}{1}
\newcounter{AlternativeCounter}
\setcounter{AlternativeCounter}{1}

\definecolor{colorUC}{RGB}{42,161,152}
\newcommand{\usecasebox}[1]{\colorbox{colorUC}{\textcolor{white}{#1}}}

\newcommand{\header}[1]{\subsection[(PU\arabic{UseCaseCounter}) #1]{\usecasebox{PU\arabic{UseCaseCounter}} #1}}

\newcommand{\desc}[1]{\subsubsection{Opis przypadku użycia}
	\noindent #1}

\newcommand{\main}[1]{\subsubsection{Podstawowy przebieg}
	\noindent #1}

\newcommand{\beginalternatives}{\subsubsection{Przebiegi alternatywne}\noindent}
\newcommand{\alternative}[2]{\noindent\begin{tabular}{|p{\textwidth}|}
		\hline
		\textbf{PU\arabic{UseCaseCounter}.\Alph{AlternativeCounter} #1}\\
		#2\\
		\hline
	\end{tabular}\vspace{0.5em}\stepcounter{AlternativeCounter}}

\newcommand{\diagram}[2]{\vspace{-0.5em}\subsubsection{Diagram przypadku użycia}\noindent
	\begin{figure}[!ht]
		\centering
		\includegraphics[width=#2\textwidth]{diagrams/#1}
	\end{figure}
	\FloatBarrier}

\newcommand{\sekwencja}[2]{\vspace{-0.5em}\subsubsection{Diagram sekwencji}\noindent
	\begin{figure}[!ht]
		\centering
		\includegraphics[width=#2\textwidth]{diagrams/#1}
	\end{figure}
	\FloatBarrier}
    
\newcommand{\klasa}[2]{\vspace{-0.5em}\subsubsection{Diagram klas}\noindent
	\begin{figure}[!ht]
		\centering
		\includegraphics[width=#2\textwidth]{diagrams/#1}
	\end{figure}
	\FloatBarrier}

\newcommand{\precond}[1]{\subsubsection{Warunki wstępne}
	\noindent #1}

\newcommand{\postcond}[1]{\subsubsection{Warunki końcowe}
	\noindent #1}

\newcommand{\references}[1]{\subsubsection{Referencje}
	\noindent #1}

\definecolor{xcolorFR}{RGB}{108,113,196}
\definecolor{xcolorNFR}{RGB}{220,50,47}
\newcommand{\xfuncreqbox}[1]{\colorbox{xcolorFR}{\textcolor{white}{#1}}}
\newcommand{\xnonfuncreqbox}[1]{\colorbox{xcolorNFR}{\textcolor{white}{#1}}}
\newcommand{\rf}[1]{\xfuncreqbox{F#1}}
\newcommand{\rnf}[1]{\xnonfuncreqbox{NF#1}}

\newcommand{\actors}[1]{\subsubsection{Aktorzy: \textnormal{#1}}}
\newcommand{\priority}[1]{\subsubsection{Priorytet: \textnormal{#1}}}


\definecolor{colorTODO}{RGB}{255,0,255}
\definecolor{colorCOMMENT}{RGB}{255,120,0}
\newcommand{\todo}{\colorbox{colorTODO}{\textcolor{white}{\textbf{TODO}}}}
\newcommand{\comment}[1]{\textcolor{colorCOMMENT}{#1}}

\newenvironment{usecase}{\setcounter{AlternativeCounter}{1}\header}{\stepcounter{UseCaseCounter}}
% Szablon do wymagan funkc. i niefunkc.
\newcounter{FuncReqCounter}
\setcounter{FuncReqCounter}{1}
\newcounter{NonFuncReqCounter}
\setcounter{NonFuncReqCounter}{1}

\definecolor{colorFR}{RGB}{108,113,196}
\definecolor{colorNFR}{RGB}{220,50,47}
\newcommand{\funcreqbox}[1]{\colorbox{colorFR}{\textcolor{white}{#1}}}
\newcommand{\nonfuncreqbox}[1]{\colorbox{colorNFR}{\textcolor{white}{#1}}}

\newcommand{\headerFR}[1]{\subsection[(F\arabic{FuncReqCounter}) #1]{\funcreqbox{F\arabic{FuncReqCounter}} #1}}
\newcommand{\headerNFR}[1]{\subsection[(NF\arabic{NonFuncReqCounter}) #1]{\nonfuncreqbox{NF\arabic{NonFuncReqCounter}} #1}}
\newcommand{\prioritytable}[2]{\begin{tabular}[ht!]{|l|c|}
		\hline
		Priorytet realizacji: & #1\\
		\hline
		Złożoność realizacji: & #2\\
		\hline
\end{tabular}}

\newcommand{\funcreq}[3]{\headerFR{#1} \prioritytable{#2}{#3} \stepcounter{FuncReqCounter}}
\newcommand{\nonfuncreq}[3]{\headerNFR{#1} \stepcounter{NonFuncReqCounter}}
% Szablon do diagramów sekwencji
\newcounter{SequenceCounter}
\setcounter{SequenceCounter}{1}

\newcommand{\headerSQ}[1]{\subsection{#1}}

\newcommand{\diagramSQ}[2]{\vspace{-0.5em}\subsubsection{}\noindent
	\begin{figure}[!ht]
		\centering
		\includegraphics[width=#2\textwidth]{diagrams/#1}
	\end{figure}
	\FloatBarrier}

\newenvironment{sequence}{\headerSQ}{\stepcounter{SequenceCounter}}
%\newcommand{\sequnece}[3]{\headerSQ{#1} \diagramSQ{#2}{#3} \stepcounter{SequenceCounter}}
% Szablon do diagramów klas
\newcounter{ClassCounter}
\setcounter{ClassCounter}{1}

\newcommand{\headerCD}[1]{\subsection{#1}}

\newcommand{\diagramCD}[2]{\vspace{-0.5em}\subsubsection{Diagram klas}\noindent
	\begin{figure}[!ht]
		\centering
		\includegraphics[width=#2\textwidth]{diagrams/#1}
	\end{figure}
	\FloatBarrier}

\newenvironment{class}{\headerSQ}{\stepcounter{ClassCounter}}
% Szablon do diagramów komponentów
\newcounter{ComponentCounter}
\setcounter{ComponentCounter}{1}

\newcommand{\headerC}[1]{\subsection{#1}}

\newcommand{\diagramC}[2]{\vspace{-0.5em}\subsubsection{}\noindent
	\begin{figure}[!ht]
		\centering
		\includegraphics[width=#2\textwidth]{diagrams/#1}
	\end{figure}
	\FloatBarrier}

% W pliku t_database.tex też jest zdefiniowane srodowisko o takiej nazwie - popraw to, bo PDF się nie chce zbudować
%\newenvironment{component}{\headerSQ}{\stepcounter{ComponentCounter}}

\begin{document}

% Title page
\begin{titlepage}
\centering
\begin{figure}
	\centering
	\includegraphics[width=0.9\textwidth]{logo.pdf}
\end{figure}
\vspace*{100pt}
\LARGE{Projekt APSI}\\
\vspace{30pt}
\textsc{\Huge{System Grapics}}\\
\vspace{120pt}
\Large{Mikołaj Gapsa}\\
\Large{Łukasz Kilaszewski}\\
\Large{Vitali Kozak}\\
\Large{Piotr Rosłaniec}\\
\vfill
\large{Warszawa, 2017}
\end{titlepage}

\thispagestyle{empty}
\tableofcontents
\cleardoublepage

\chapter{Wprowadzenie}
\section{Opis projektu}
Projekt Grapics realizuje system informatyczny służący do zarządzania procesami biznesowymi w firmie realizującej zamówienia B2B na dobra intelektualne, w szczególności grafiki i fotografie.

\section{Procesy biznesowe}
Ze względu na zróżnicowane zadania jakie realizuje wykonawca (odbiorca systemu) proces realizacja zadania nie jest dobrze ustrukturyzowany. Utrudnia to monitorowania i nadzorowanie zadań, a także ich synchronizację między sobą w przypadkach, kiedy istnieją między nimi zależności.

Zmienność wymagań i ich niedoprecyzowanie sprawia, że konieczna jest ciągła komunikacja z klientem i pozyskiwanie informacji zwrotnej po wprowadzeniu znaczących zmian. Efektem każdej zmiany jest wydanie artefaktu w pewnej wersji. Artefakty wymagają nie tylko akceptacji klienta, ale również konkretnych osób z zespołu i przełożonego. 

Śledzenie i wersjonowanie artefaktów staje się konieczne w projektach, które przechodzą wiele iteracji nierzadko przy udziale podwykonawców. Artefakty mogą się znajdować w wielu różnych lokalizacjach, po stronie klienta, wykonawcy lub podwykonawców.

\section{Opis wprowadzonych pojęć}
W ramach opisu projektu wprowadzono następujące pojęcia:
\begin{itemize}
	\item \textit{Zgłoszenie} -- jednostka pracy o znaczeniu biznesowym,
	\item \textit{Artefakt} -- produkt realizacji \textit{Zgłoszenia},
	\item \textit{Projekt} -- zbiór \textit{Zgłoszeń},
	\item \textit{Repozytorium} -- lokalizacja przechowywująca \textit{Artefakty}.
\end{itemize}

\section{Aktorzy}
W systemie zidentyfikowano następujących aktorów:
\begin{itemize}
	\item \textit{Administrator} -- jednostka pracy o znaczeniu biznesowym,
	\item \textit{Użytkownik} -- podstawowy aktor, jego rola biznesowa jest określona za pomocą zestawu uprawnień dla poszczególnych projektów,
	\item \textit{Twórca projektu} (\texttt{<<extends>>} \textit{Użytkownik}) -- użytkownik sprawujący odpowiedzialność nad projektem,
	\item \textit{Twórca zgłoszenia} (\texttt{<<extends>>} \textit{Użytkownik}) -- użytkownik sprawujący odpowiedzialność nad zgłoszeniem,
	\item \textit{Osoba Odpowiedzialna} (\texttt{<<extends>>} \textit{Użytkownik}) -- użytkownik oddelegowany do zrealizowania zgłoszenia,
	\item \textit{Gość} -- osoba nie posiadająca dostępu do systemu.
\end{itemize}

\section{Moduły}
System podzielony jest na następujące moduły:
\begin{itemize}
	\item Moduł uwierzytelniania
	\item Moduł administracyjny
	\item Moduł główny (obsługi projektów)	
\end{itemize}
\subsection{Moduł uwierzytelniania}
Moduł odpowiedzialny za uwierztyelnianie użytkowników. Jego główne cele to umożliwienie gościom zakładanie kont oraz pozwolić użytkownikom zalogować się.
\subsection{Moduł administracyjny}
Moduł odpowiedzialny za wszelkie prace wykonywane przez Administratora Systemu. Uwzględniono tutaj wszystkie możliwości jakie można znaleźć w panelu Administratora.
\subsection{Moduł główny}
Opis \ldots \todo

\chapter{Specyfikacja sprzętu i oprogramowania podstawowego}
\todo
\chapter{Specyfikacja technologii}
\todo
\chapter{Przypadki użycia}
\section{Moduł uwierzytelniania}
\begin{usecase}{Uwierzytelnianie użytkownika}
	% Opis przypadku użycia
	\desc{Ten ogólny przypadek użycia opisuje proces uwierzytelniania użytkowników w systemie Grapics.}
    % Priorytet
	\priority{10}
	% Aktorzy
	\actors{Użytkownik, Gość}
	% Przebieg główny
	\main{Ten przypadek użycia zaczyna się gdy Użytkownik chcę korzystać z Systemu Grapics.}
	\begin{enumerate}
		\item Logowanie
        \item Rejestracja
	\end{enumerate}
    \beginalternatives
	\alternative{Użytkownik nie jest zarejestrowany}{Jeśli użytkownik nie jest jeszcze zarejstrowany, nie może on się zalogować. Jeśli nie przejdzie procesu rejestracji, nie będzie on mógł pomyślnie przejśćprocesu uwierzytelniania co kończy ten przypadek użycia. }
	%\alternative{Dodatkowa pozycja}{W ten sposób można dodawać kolejne przebiegi alternatywne \ldots}
	% Diagram
	\diagram{uwierzytelnianie.pdf}{1} % drugi argument to szerokość obrazu (1.0 = szerokość strony)
	% Pre-conditions
	\precond{Brak. }
	% Post-conditions
	\postcond{Jeśli przypadek użycia został pomyślnie zrealizowany, aktor jest uewierzytelniony w Systemie.}
    % Referencje
	\references{\rf{1} \rf{2} \rf{4}}
\end{usecase}

\begin{usecase}{Logowanie}
\label{logowanie}
	% Opis przypadku użycia
	\desc{Ten przypadek użycia opisuje jak użytkownik loguje się do systemu Grapics.}
    % Priorytet
	\priority{10}
	% Aktorzy
	\actors{Użytkownik}
	% Przebieg główny
	\main{Ten przypadek użycia zaczyna się gdy użytkownik próbuje zalogować się do systemu.}
	\begin{enumerate}
		\item System prosi Użytkownika o zalogowanie
		\item Użytkownik podaje swoją nazwę użytkownika oraz hasło
		\item System potwierdza poprawność danych i loguje użytkownika do systemu
	\end{enumerate}
	% Przebiegi alternatywne
	\beginalternatives
	\alternative{Nieprawidłowa nazwa użytkownika/hasło}{Jeśli użytkownik podał złą nazwą użytkownika i/lub hasło System wyświetla odpowiedni komunikat błędu. Aktor w tym momencie może zdecydować by wrócić do początku przebiegu podstawowego lub anulowa logowanie, co kończy ten przypadek użycia}
	%\alternative{Dodatkowa pozycja}{W ten sposób można dodawać kolejne przebiegi alternatywne \ldots}
	% Diagram
	\diagram{logowanie.pdf}{1} % drugi argument to szerokość obrazu (1.0 = szerokość strony)
    % Diagram sekwencji
	\sekwencja{logowanie_sekwencja.pdf}{1}
	% Pre-conditions
	\precond{Użytkownik jest zarejestrowany w systemie.}
	% Post-conditions
	\postcond{Jeśli przypadek użycia został pomyślnie zrealizowany, aktor jest zalogowany do systemu. W przeciwnej sytuacji stan systemu pozostaje niezmieniony.}
    % Referencje
	\references{\rf{2}}
\end{usecase}

\begin{usecase}{Rejestracja}
	% Opis przypadku użycia
	\desc{Ten przypadek użycia opisuje proces zakładania konta w systemu Grapics.}
    % Priorytet
	\priority{10}
	% Aktorzy
	\actors{Administrator, Gość}
	% Przebieg główny
	\main{Ten przypadek użycia zaczyna się gdy użytkownik próbuje zarejestrować się w systemie.}
	\begin{enumerate}
		\item System prosi Użytkownika o wypełnienie formularza rejestracyjnego
		\item Użytkownik wypełnia formularz
		\item System potwierdza poprawność danych
        \item Administrator zatwierdza konto nowo-zarejestrowanego użytkownika
	\end{enumerate}
	% Przebiegi alternatywne
	\beginalternatives
	\alternative{Nieprawidłowa nazwa użytkownika/hasło}{Jeśli użytkownik podał znazwę użytkownika, która już jest używana w systemie lub/i hasło, które nie spełnia wymogów systemowych, System wyświetla odpowiedni komunikat błędu. Aktor w tym momencie może zdecydować by wrócić do początku przebiegu podstawowego lub anulować rejestrację, co kończy ten przypadek użycia}
    \alternative{Administrator nie zatwierdzi konta}{Administrator zatwierdza tylko konta, których dane pokrywają się z wewnętrzną listą klientów. Jeśli nieuprawniona osoba próbowała założyć konto, nie zostanie ono zatwierdzone co kończy ten przypadek użycia}
	%\alternative{Dodatkowa pozycja}{W ten sposób można dodawać kolejne przebiegi alternatywne \ldots}
	% Diagram
	\diagram{rejestracja.pdf}{1} % drugi argument to szerokość obrazu (1.0 = szerokość strony)
    % Diagram sekwencji
	\sekwencja{rejestracja_sekwencja.pdf}{1}
	% Pre-conditions
	\precond{Brak.}
	% Post-conditions
	\postcond{Jeśli przypadek użycia został pomyślnie zrealizowany, aktor jest zarejestrowany w Systemie i może przytąpić do logowania się. W przeciwnej sytuacji stan systemu pozostaje niezmieniony.}
    % Referencje
	\references{\rf{1} \rf{4} \rnf{1} \rnf{2}}
\end{usecase}


\section{Moduł administracyjny}
\begin{usecase}{Zarządzanie Systemem}
	% Opis przypadku użycia
	\desc{Ten ogólny przypadek użycia opisuje jak Administrator może zarządzać Systemem.}
    % Priorytet
	\priority{9}
	% Aktorzy
	\actors{Administrator}
	% Przebieg główny
	\main{Ten przypadek użycia zaczyna się gdy Administrator wejdzie do panelu Administratora}
	\begin{enumerate}
		\item Administrator robi backup danych
		\item Administrator wykonuje przywracanie systemu
		\item Administrator przegląda logi systemu
        \item Administrator przeprwowadza rekonfigurację systemu
        \item Administrator zarządza kontem Użytkownika
	\end{enumerate}
	% Przebiegi alternatywne
	\beginalternatives
	\alternative{Administrator próbował wykonać przywracanie systemu w niedozwolonym momencie}{ System nie pozwoli Administratorowi wykonać takiej opracji w innym momencie niż ten ściśle określony w kontrakcie}
	%\alternative{Dodatkowa pozycja}{W ten sposób można dodawać kolejne przebiegi alternatywne \ldots}
	% Diagram
	\diagram{zarzadzanie_systemem.pdf}{1} % drugi argument to szerokość obrazu (1.0 = szerokość strony)
	% Pre-conditions
	\precond{Administrator jest zalogowany w systemie. }
	% Post-conditions
	\postcond{Jeśli przypadek użycia został pomyślnie zrealizowany, stan danego konta użytkownika został zmieniony.}
    % Referencje
	\references{\rf{5} \rf{6} \rf{7} \rf{8} \rf{9} \rf{10} \rnf{2} \rnf{4} \rnf{5}}
\end{usecase}

\begin{usecase}{Zarządzanie kontem Użytkownika}
	% Opis przypadku użycia
	\desc{Ten ogólny przypadek użycia opisuje jak Administrator zarządza kontem użytkownika.}
    % Priorytet
	\priority{9}
	% Aktorzy
	\actors{Administrator}
	% Przebieg główny
	\main{Ten przypadek użycia zaczyna się gdy Administrator wejdzie do panelu Administratora}
	\begin{enumerate}
		\item Administrator Zatwierdza konto nowo-zarejestrowanego Użytkownika
		\item Administrator zmienia uprawnienia użytkownika
		\item Administrator usuwa konto użytkownika
	\end{enumerate}
	% Przebiegi alternatywne
	\beginalternatives
	\alternative{Administrator próbował zarządać kontem innego Administratora}{ Jeśli Administrator próbował zarządzać kontem innego Administratora, System odrzuca zmiany co kończy ten przypadek użycia}
	%\alternative{Dodatkowa pozycja}{W ten sposób można dodawać kolejne przebiegi alternatywne \ldots}
	% Diagram
	\diagram{zarzadzanie_kontem.pdf}{1} % drugi argument to szerokość obrazu (1.0 = szerokość strony)
	% Pre-conditions
	\precond{Administrator jest zalogowany w Systemie. }
	% Post-conditions
	\postcond{Jeśli przypadek użycia został pomyślnie zrealizowany, stan danego konta użytkownika został zmieniony.}
    % Referencje
	\references{\rf{5} \rf{6} \rf{7}}
\end{usecase}

\begin{usecase}{Zatwierdzanie konta}
	% Opis przypadku użycia
	\desc{Ten przypadek użycia opisuje proces zatwierdzania kont nowo-zarejestrowanych użytkowników w systemie Grapics.}
    % Priorytet
	\priority{10}
	% Aktorzy
	\actors{Administrator}
	% Przebieg główny
	\main{Ten przypadek użycia zaczyna się gdy Administrator otrzymał powiadomienie o nowo-utworzonych kontach się w systemie.}
	\begin{enumerate}
		\item System wyświetla nowych użytkowników w panelu Administratorskim
        \item Administrator wybiera jedno z kont
		\item Administrator weryfikuje i zatwierdza wybrane konto
	\end{enumerate}
	% Przebiegi alternatywne
	\beginalternatives
	\alternative{Osoba nieupowazniona próbowała założyć konto}{Jeśli osoba nieupoważniona, to znaczy taka, która nie jest klientem, próbowała założyć konto w systemie, konto zostaje odrzucone przez administratora}
	%\alternative{Dodatkowa pozycja}{W ten sposób można dodawać kolejne przebiegi alternatywne \ldots}
	% Diagram
	\diagram{zatwierdzanie_konta.pdf}{1} % drugi argument to szerokość obrazu (1.0 = szerokość strony)
    % Diagram sekwencji
	\sekwencja{zatwierdzenie_konta_sekwencja.pdf}{1}
	% Diagram klas
	\klasa{zatwierdzenie_konta_klas.pdf}{1}
	% Pre-conditions
	\precond{Administrator jest zalogowany w Systemie oraz istnieją niezatwierdzone konta w Systemie.}
	% Post-conditions
	\postcond{Jeśli przypadek użycia został pomyślnie zrealizowany, konto jest potwierdzone przez Administratora w Systemie i może on posłużyć do zalogowania się. W przeciwnej sytuacji stan systemu pozostaje niezmieniony.}
    % Referencje
	\references{\rf{5}}
\end{usecase}

\begin{usecase}{Usuwanie konta}
	% Opis przypadku użycia
	\desc{Ten przypadek użycia opisuje proces usuwania kont użytkowników w systemie Grapics.}
    % Priorytet
	\priority{8}
	% Aktorzy
	\actors{Administrator}
	% Przebieg główny
	\main{Ten przypadek użycia zaczyna się gdy Administrator otrzymał powiadomienie o wygaścięciu ważności konta w Systemie Grapics.}
	\begin{enumerate}
		\item System wyświetla użytkowników w panelu Administratorskim
        \item Administrator wybiera jedno z kont
		\item Administrator weryfikuje i usuwa wybrane konto
	\end{enumerate}
	% Przebiegi alternatywne
	%\alternative{Dodatkowa pozycja}{W ten sposób można dodawać kolejne przebiegi alternatywne \ldots}
	% Diagram
	\diagram{usuwanie_konta.pdf}{1} % drugi argument to szerokość obrazu (1.0 = szerokość strony)
    % Diagram sekwencji
	\sekwencja{usuwanie_konta_sekwencja.pdf}{1}
	% Diagram klas
	\klasa{usuwanie_konta_zmiana_uprawnien_klas.pdf}{1}
	% Pre-conditions
	\precond{Administrator jest zalogowany w Systemie oraz dostał powiadomienie o wygaśnięciu wazności konta użytkownika.}
	% Post-conditions
	\postcond{Jeśli przypadek użycia został pomyślnie zrealizowany, konto jest usunięte przez Administratora w Systemie i nie może ono posłużyć do zalogowania się. W przeciwnej sytuacji stan systemu pozostaje niezmieniony.}
    % Referencje
	\references{\rf{7}}
\end{usecase}

\begin{usecase}{Zmiana uprawnień}
	% Opis przypadku użycia
	\desc{Ten przypadek użycia opisuje proces zmieniania uprawnież kont użytkowników w systemie Grapics.}
    % Priorytet
	\priority{8}
	% Aktorzy
	\actors{Administrator}
	% Przebieg główny
	\main{Ten przypadek użycia zaczyna się gdy Administrator otrzymał powiadomienie o potrzebie zmiany uprawnień konta w Systemie Grapics.}
	\begin{enumerate}
		\item System wyświetla użytkowników w panelu Administratorskim
        \item Administrator wybiera jedno z kont
		\item Administrator weryfikuje i zmienia uprawnienia wybranengo konta
	\end{enumerate}
	% Przebiegi alternatywne
	%\alternative{Dodatkowa pozycja}{W ten sposób można dodawać kolejne przebiegi alternatywne \ldots}
	% Diagram
	\diagram{zmiana_uprawnien.pdf}{1} % drugi argument to szerokość obrazu (1.0 = szerokość strony)
    % Diagram sekwencji
	\sekwencja{zmiana_uprawnien_sekwencja.pdf}{1}
	% Diagram klas
	\klasa{usuwanie_konta_zmiana_uprawnien_klas.pdf}{1}
	% Pre-conditions
	\precond{Administrator jest zalogowany w Systemie oraz dostał powiadomienie o potrzebie zmiany uprawnień danego konta użytkownika.}
	% Post-conditions
	\postcond{Jeśli przypadek użycia został pomyślnie zrealizowany, uprawnienia danego konta zostały zmienione. W przeciwnej sytuacji stan systemu pozostaje niezmieniony.}
    % Referencje
	\references{\rf{6}}
\end{usecase}

\begin{usecase}{Backup danych}
	% Opis przypadku użycia
	\desc{Ten przypadek użycia opisuje proces robienia backupu danych przez Administratora.}
    % Priorytet
	\priority{10}
	% Aktorzy
	\actors{Administrator}
	% Przebieg główny
	\main{Ten przypadek użycia zaczyna się gdy Administrator otrzymał powiadomienie o potrzebie zrobienia backupu danych w Systemie Grapics.}
	\begin{enumerate}
		\item Administrator wybiera opcję "Backup danych" w panelu Administratorskim
        \item Administrator zaznacza odpowiednie opcje
		\item Administrator czeka aż proces dobiegnie końca
	\end{enumerate}
	% Przebiegi alternatywne
	%\alternative{Dodatkowa pozycja}{W ten sposób można dodawać kolejne przebiegi alternatywne \ldots}
	% Diagram
	\diagram{backup_danych.pdf}{1} % drugi argument to szerokość obrazu (1.0 = szerokość strony)
    % Diagram sekwencji
	\sekwencja{backup_sekwencja.pdf}{1}
	% Diagram klas
	\klasa{backup_klas.pdf}{1}
	% Pre-conditions
	\precond{Administrator jest zalogowany w Systemie oraz dostał powiadomienie o potrzebie wykonania Backupu.}
	% Post-conditions
	\postcond{Jeśli przypadek użycia został pomyślnie zrealizowany, backup Systemu został wykonany. W przeciwnej sytuacji stan systemu pozostaje niezmieniony.}
    % Referencje
	\references{\rf{8} \rnf{2}}
\end{usecase}

\begin{usecase}{Przywracanie Systemu}
	% Opis przypadku użycia
	\desc{Ten przypadek użycia opisuje proces przywracania Systemu przez Administratora.}
    % Priorytet
	\priority{10}
	% Aktorzy
	\actors{Administrator}
	% Przebieg główny
	\main{Ten przypadek użycia zaczyna się gdy Administrator otrzymał powiadomienie o potrzebie przywrócenia Systemu Grapics.}
	\begin{enumerate}
		\item Administrator odnajduje odpowiednią wersję Systemu, którą należy przywrócić
        \item Administrator wybiera opcję "Przywracanie Systemu" z panelu Administratora
        \item Administrator zaznacza odpowiednie opcje
		\item Administrator czeka aż proces dobiegnie końca
	\end{enumerate}
	% Przebiegi alternatywne
    \beginalternatives
	\alternative{Administrator próbował wykonać przywracanie systemu w niedozwolonym momencie}{ System nie pozwoli Administratorowi wykonać takiej opracji w innym momencie niż ten ściśle określony w kontrakcie}
	%\alternative{Dodatkowa pozycja}{W ten sposób można dodawać kolejne przebiegi alternatywne \ldots}
	% Diagram
	\diagram{przywracanie_systemu.pdf}{1} % drugi argument to szerokość obrazu (1.0 = szerokość strony)
    % Diagram sekwencji
	\sekwencja{przywracanie_systemu_sekwencja.pdf}{1}
	% Diagram klas
	\klasa{przywracanie_systemu_klas.pdf}{1}
	% Pre-conditions
	\precond{Administrator jest zalogowany w Systemie oraz dostał powiadomienie o potrzebie przywrócenia Systemu.}
	% Post-conditions
	\postcond{Jeśli przypadek użycia został pomyślnie zrealizowany, odpowiednia wersja Systemu została przywrócona. W przeciwnej sytuacji stan systemu pozostaje niezmieniony.}
    % Referencje
	\references{\rf{10}}
\end{usecase}

\begin{usecase}{Przeglądanie logów Systemu}
	% Opis przypadku użycia
	\desc{Ten przypadek użycia opisuje przeglądania logów systemu przez Administratora.}
    % Priorytet
	\priority{7}
	% Aktorzy
	\actors{Administrator}
	% Przebieg główny
	\main{Ten przypadek użycia zaczyna się gdy Administrator zamierza przejrzeć logi w Systemie Grapics.}
	\begin{enumerate}
		\item Administrator wybiera opcję "Logi" w panelu Administratora
        \item Administrator zaznacza odpowiednie opcje
		\item Administrator czeka aż zbieranie logów dobiegnie końca
        \item Administrator otrzymuje logi w formacie .txt
	\end{enumerate}
	% Przebiegi alternatywne
	%\alternative{Dodatkowa pozycja}{W ten sposób można dodawać kolejne przebiegi alternatywne \ldots}
	% Diagram
	\diagram{przegladanie_logow.pdf}{1} % drugi argument to szerokość obrazu (1.0 = szerokość strony)
    % Diagram sekwencji
	\sekwencja{przegladanie_logow_sekwencja.pdf}{1}
	% Diagram klas
	\klasa{przegladanie_logow_klas.pdf}{1}
	% Pre-conditions
	\precond{Administrator jest zalogowany w Systemie.}
	% Post-conditions
	\postcond{Jeśli przypadek użycia został pomyślnie zrealizowany, logi Systemu zostały wygenerowane do pliku .txt. W przeciwnej sytuacji stan systemu pozostaje niezmieniony.}
    % Referencje
	\references{\rf{9}}
\end{usecase}

\begin{usecase}{Rekonfiguracja Systemu}
	% Opis przypadku użycia
	\desc{Ten przypadek użycia opisuje proces rekonfiguracji Systemu przez Administratora.}
    % Priorytet
	\priority{9}
	% Aktorzy
	\actors{Administrator}
	% Przebieg główny
	\main{Ten przypadek użycia zaczyna się gdy Administrator otrzymał powiadomienie o potrzebie rekonfiguracji Systemu Grapics.}
	\begin{enumerate}
		\item Administrator wybiera opcję "Rekonfiguracja Systemu" w panelu Administratorskim
        \item Administrator zaznacza odpowiednie opcje
		\item Administrator czeka aż proces dobiegnie końca
	\end{enumerate}
	% Przebiegi alternatywne
    \beginalternatives
	\alternative{Administrator próbował wykonać Rekonfiguracje Systemu w niedozwolonym momencie}{ System nie pozwoli Administratorowi wykonać takiej opracji w innym momencie niż ten ściśle określony w kontrakcie}
	%\alternative{Dodatkowa pozycja}{W ten sposób można dodawać kolejne przebiegi alternatywne \ldots}
	% Diagram
	\diagram{rekonfiguracja_systemu.pdf}{1} % drugi argument to szerokość obrazu (1.0 = szerokość strony)
    % Diagram sekwencji
	\sekwencja{rekonfiguracja_sekwencja.pdf}{1}
	% Diagram klas
	\klasa{rekonfiguracja_klas.pdf}{1}
	% Pre-conditions
	\precond{Administrator jest zalogowany w Systemie oraz dostał powiadomienie o potrzebie rekonfiguracji Systemu.}
	% Post-conditions
	\postcond{Jeśli przypadek użycia został pomyślnie zrealizowany, rekonfiguracja Systemu została wykonana. W przeciwnej sytuacji stan systemu pozostaje niezmieniony.}
    % Referencje
	\references{\rf{10}}
\end{usecase}
\section{Moduł główny}
\begin{usecase}{Przeglądanie Projektów}
\label{przegladanie_projektow}
	% Opis przypadku użycia
	\desc{Użytkownik przegląda Projekty, w których funkcjonuje w roli Członka Projektu}
	% Priorytet
	\priority{10}
	% Aktorzy
	\actors{Członek Projektu}
	% Przebieg główny
	\begin{enumerate}
      \item Użytkownik wchodzi w sekcję "moje projekty"
      \item Użytkownik zaznacza opcję filtrowania projektów, w których jest członkiem
      \item System wyświetla listę projektów, w których Użytkownik funkcjonuje w roli Członka Projektu.
	\end{enumerate}
	% Diagram
	\diagram{przegladanie_projektow.pdf}{0.9} % drugi argument to szerokość obrazu (1.0 = szerokość strony)
	% Pre-conditions
	% Diagram sekwencji
	\parindent=0cm
	\sekwencja{przegladanie_projektow_sekwencja_fit.pdf}{0.9}
	% Diagram klas
	%\klasa{przegladanie_logow_klas.pdf}{1}
	\precond{Użytkownik jest zalogowany (\ref{logowanie})}
	\references{\rf{13}}
\end{usecase}


\begin{usecase}{Tworzenie Projektu}
\label{tworzenie_projektu}
	% Opis przypadku użycia
	\desc{Użytkownik tworzy nowy Projekt}
	% Priorytet
	\priority{10}
	% Aktorzy
	\actors{Członek Projektu, Twórca Projektu}
	% Przebieg główny
	\begin{enumerate}
    \item Użytkownik przegląda istniejące Projekty (\ref{przegladanie_projektow})
    \item Użytkownik wybiera opcję  “Utwórz nowy projekt”
    \item Użytkownik wypełnia formularz
    \item Użytkownik zatwierdza utworzenie nowego projektu.
	\end{enumerate}
	% Diagram
	\diagram{tworzenie_projektu.pdf}{0.9} % drugi argument to szerokość obrazu (1.0 = szerokość strony)
	% Pre-conditions
	\precond{Użytkownik jest zalogowany}
	% Post-conditions
	\postcond{Został utworzony nowy Projekt, a Użytkownik, który go stworzył, zostaje Twórcą tego Projektu}
	\references{\rf{11}}
\end{usecase}


\begin{usecase}{Edycja Projektu}
	% Opis przypadku użycia
	\desc{Twórca Projektu edytuje swój Projekt}
	% Priorytet
	\priority{10}
	% Aktorzy
	\actors{Członek Projektu, Twórca Projektu}
	% Przebieg główny
	\begin{enumerate}
    \item Twórca Projektu przegląda istniejące Projekty (\ref{przegladanie_projektow})
    \item Twórca Projektu wybiera Projekt
    \item Twórca Projektu dokonuje zmian w formularzu
    \item Twórca Projektu zatwierdza zmiany
	\end{enumerate}
	% Diagram
	\diagram{edycja_projektu.pdf}{0.9} % drugi argument to szerokość obrazu (1.0 = szerokość strony)
	% Diagram sekwencji
	%\sekwencja{edycja_projektu_sekwencja_fit.pdf}{0.9}
	\parindent=0cm
	\textbf{Diagram sekwencji}
	\diagramSQ{edycja_projektu_sekwencja_fit.pdf}{0.9}
	% Pre-conditions
	\precond{Użytkownik jest zalogowany (\ref{logowanie}) i jest Twórcą edytowanego Projektu (\ref{tworzenie_projektu})}
	% Post-conditions
	\postcond{Właściwości Projektu zostały zmienione}
	\references{\rf{19}}
\end{usecase}


\begin{usecase}{Zamknięcie projektu}
	% Opis przypadku użycia
	\desc{Twórca Projektu zgłasza do zamknięcia swój Projekt}
	% Priorytet
	\priority{10}
	% Aktorzy
	\actors{Twórca Projektu}
	% Przebieg główny
	\begin{enumerate}
	\item Twórca Projektu przegląda istniejące Projekty (\ref{przegladanie_projektow})
    \item Twórca Projektu wybiera Projekt
    \item Twórca Projektu wybiera Zgłoszenie
    \item Twórca Projektu wydaje polecenie zamknięcia projektu
	\end{enumerate}
	% Diagram
	\diagram{zamkniecie_projektu.pdf}{0.9} % drugi argument to szerokość obrazu (1.0 = szerokość strony)
	% Diagram sekwencji
	%\sekwencja{zamkniecie_projektu_sekwencja_fit.pdf}{0.9}
	\parindent=0cm
	\textbf{Diagram sekwencji}
	\diagramSQ{zamkniecie_projektu_sekwencja_fit.pdf}{0.9}
	% Pre-conditions
	\precond{Użytkownik jest zalogowany (\ref{logowanie}) i jest Twórcą zamykanego Projektu (\ref{tworzenie_projektu})}
	% Post-conditions
	\postcond{Projekt został zamknięty, Użytkownik nie jest już Twórcą Projektu}
	\references{\rf{20}}
\end{usecase}


\begin{usecase}{Tworzenie Zgłoszenia}
\label{tworzenie_zgloszenia}
	% Opis przypadku użycia
	\desc{Opisuje proces, w którym Użytkownik będący Członkiem Projektu dodaje nowe Zgłoszenie do Projektu}
	% Priorytet
	\priority{10}
	% Aktorzy
	\actors{Członek Projektu}
	% Przebieg główny
	\begin{enumerate}
      \item Członek Projektu przegląda istniejące Projekty (\ref{przegladanie_projektow})
      \item Członek Projektu wybiera odpowiedni Projekt
      \item Członek Projektu uzupełnia szczegóły nowego Zgłoszenia
      \item Członek Projektu zatwierdza powstanie nowego Zgłoszenia
	\end{enumerate}
	% Przebiegi alternatywne
	\beginalternatives
	\alternative{Wybranie Osoby Odpowiedzialnej za Zgłoszenie}{Członek Projektu wybiera Osobę Odpowiedzialną za utworzone zgłoszenie spośród Członków Projektu (\ref{wyznaczenie_osoby_odpowiedzialnej})}
	% Diagram
	\diagram{tworzenie_zgloszenia.pdf}{0.9} % drugi argument to szerokość obrazu (1.0 = szerokość strony)
	% Diagram sekwencji
	%\sekwencja{tworzenie_zgloszenia_sekwencja_fit.pdf}{0.9}
	\parindent=0cm
	\textbf{Diagram sekwencji}
	\diagramSQ{tworzenie_zgloszenia_sekwencja_fit.pdf}{0.9}
	% Pre-conditions
	\precond{Użytkownik (Członek Projektu) jest zalogowany (\ref{logowanie}) i jest członkiem przynajmniej jednego Projektu}
	% Post-conditions
	\postcond{Powstało nowe Zgłoszenie w Projekcie, Członek Projektu staje się Twórca Zgłoszenia dla utworzonego Zgłoszenia}
	\references{\rf{14}}
\end{usecase}


\begin{usecase}{Zamknięcie Zgłoszenia}
	% Opis przypadku użycia
	\desc{Twórca Zgłoszenia zamyka swoje Zgłoszenie}
	% Priorytet
	\priority{10}
	% Aktorzy
	\actors{Twórca Zgłoszenia}
	\begin{enumerate}
    \item Twórca Zgłoszenia przegląda istniejące Projekty (\ref{przegladanie_projektow})
    \item Twórca Zgłoszenia wybiera Projekt
    \item Twórca Zgłoszenia wybiera Zgłoszenie
    \item Twórca Zgłoszenia wydaje polecenie zamknięcia zgłoszenia
    \end{enumerate}
	% Diagram
	\diagram{zamkniecie_zgloszenia.pdf}{0.9} % drugi argument to szerokość obrazu (1.0 = szerokość strony)
	% Diagram sekwencji
	%\sekwencja{zamkniecie_zgloszenia_sekwencja_fit.pdf}{0.9}
	\parindent=0cm
	\textbf{Diagram sekwencji}
	\diagramSQ{zamkniecie_zgloszenia_sekwencja_fit.pdf}{0.9}
	% Pre-conditions
    \precond{Użytkownik jest zalogowany (\ref{logowanie}) i jest Twórcą zamykanego Zgłoszenia (\ref{tworzenie_zgloszenia})}
	% Post-conditions
	\postcond{Zgłoszenie zostało zamknięte, Użytkownik nie jest już Twórcą Zgłoszenia}
	\references{\rf{21}}
\end{usecase}


\begin{usecase}{Zmiana widoczności Zgłoszenia}
	% Opis przypadku użycia
	\desc{Twórca Zgłoszenia zmienia widoczność swojego Zgłoszenia}
	% Priorytet
	\priority{5}
	% Aktorzy
	\actors{Twórca Zgłoszenia}
    % Przebieg główny
	\begin{enumerate}
    \item Twórca Zgłoszenia przegląda istniejące Projekty (\ref{przegladanie_projektow})
    \item Twórca Zgłoszenia wybiera Projekt
    \item Twórca Zgłoszenia wybiera Zgłoszenie
    \item Twórca Zgłoszenia wydaje polecenia ukrycia/odkrycia Zgłoszenia
	\end{enumerate}
	% Diagram
	\diagram{zmiana_widocznosci_zgloszenia.pdf}{0.9} % drugi argument to szerokość obrazu (1.0 = szerokość strony)
	% Diagram sekwencji
	%\sekwencja{zmiana_widocznosci_zgloszenia_sekwencja_fit.pdf}{0.9}
	\parindent=0cm
	\textbf{Diagram sekwencji}
	\diagramSQ{zmiana_widocznosci_zgloszenia_sekwencja_fit.pdf}{0.9}
	% Pre-conditions
	\precond{Użytkownik jest zalogowany (\ref{logowanie}) i jest Twórcą edytowanego Zgłoszenia (\ref{tworzenie_zgloszenia})}
	% Post-conditions
	\postcond{Widoczność Zgłoszenia została zmieniona}
		\references{\rf{22}}
\end{usecase}


\begin{usecase}{Edycja Zgłoszenia}
\label{edycja_zgloszenia}
	% Opis przypadku użycia
	\desc{Twórca Zgłoszenia edytuje podstawowe informacje o Zgłoszeniu}
	% Priorytet
	\priority{9}
	% Aktorzy
	\actors{Twórca Zgłoszenia}
    % Przebieg główny
	\begin{enumerate}
	\item Twórca Zgłoszenia przegląda istniejące Projekty (\ref{przegladanie_projektow})
    \item Twórca Zgłoszenia wybiera Projekt
    \item Twórca Zgłoszenia wybiera Zgłoszenie
    \item Twórca Zgłoszenia wydaje polecenie edycji Zgłoszenia
    \item Twórca Zgłoszenia edytuje formularz
    \item Twórca Zgłoszenia zatwierdza zmiany
	\end{enumerate}
	% Diagram
	\diagram{edycja_zgloszenia.pdf}{0.9} % drugi argument to szerokość obrazu (1.0 = szerokość strony)
	% Diagram sekwencji
	%\sekwencja{edycja_zgloszenia_sekwencja_fit.pdf}{0.9}
	\parindent=0cm
	\textbf{Diagram sekwencji}
	\diagramSQ{edycja_zgloszenia_sekwencja_fit.pdf}{0.9}
	% Pre-conditions
	\precond{Użytkownik jest zalogowany (\ref{logowanie}), należy do przynajmniej jednego Projektu i jest Twórcą przynajmniej jednego Zgłoszenia (\ref{tworzenie_zgloszenia})}
	% Post-conditions
	\postcond{Właściwości Zgłoszenia zostały zmienione}
	\references{\rf{23}}
\end{usecase}


\begin{usecase}{Wyznaczenie Osoby Odpowiedzialnej}
\label{wyznaczenie_osoby_odpowiedzialnej}
	% Opis przypadku użycia
	\desc{Twórca Zgłoszenia wyznacza Użytkownika jako Odbiorcę Zgłoszenia dla danego Zgłoszenia}
	% Priorytet
	\priority{9}
	% Aktorzy
	\actors{Twórca Zgłoszenia, Osoba Odpowiedzialna, Użytkownik}
	% Przebieg główny
	\begin{enumerate}
    \item Twórca Zgłoszenia przegląda istniejące Projekty (\ref{przegladanie_projektow})
  	\item Twórca Zgłoszenia wybiera Projekt
  	\item Twórca Zgłoszenia wybiera Zgłoszenie
  	\item Twórca Zgłoszenia wydaje polecenie edycji Zgłoszenia (\ref{edycja_zgloszenia})
  	\item Twórca Zgłoszenia wybiera z listy Użytkowników Osobę Odpowiedzialną
  	\item Twórca Zgłoszenia zatwierdza zmiany
	\end{enumerate}
	% Diagram
	\diagram{wyznaczenie_osoby_odpowiedzialnej.pdf}{0.9} % drugi argument to szerokość obrazu (1.0 = szerokość strony)
	% Pre-conditions
	\precond{Użytkownik jest zalogowany (\ref{logowanie}), należy do przynajmniej jednego Projektu i jest Twórcą przynajmniej jednego Zgłoszenia (\ref{tworzenie_zgloszenia})}
	% Post-conditions
	\postcond{Do zgłoszenia została przypisana osoba odpowiedzialna}
	\references{\rf{24}}
\end{usecase}


\begin{usecase}{Edycja Grafu}
	% Opis przypadku użycia
	\desc{Twórca Projektu planuje kolejne etapy Projektu w formie Zgłoszeń, przechowywanych w postaci Grafu}
	% Priorytet
	\priority{10}
	% Aktorzy
	\actors{Twórca Projektu}
	% Przebieg główny
	\begin{enumerate}
    \item Twórca Projektu przegląda istniejące Projekty (\ref{przegladanie_projektow})
    \item Twórca Projektu wybiera Projekt
    \item Twórca Projektu wybiera istniejące Zgłoszenie
    \item Twórca Projektu wyznacza rodzica Zgłoszenia lub usuwa istniejące dzieci
	\end{enumerate}
	% Przebiegi alternatywne
	% Diagram
	\diagram{tworzenie_grafu.pdf}{0.9} % drugi argument to szerokość obrazu (1.0 = szerokość strony)
	% Pre-conditions
	\precond{Użytkownik jest zalogowany (\ref{logowanie})  i jest Twórcą przynajmniej jednego Projektu (\ref{tworzenie_projektu})}
	% Post-conditions
	\postcond{Struktura Grafu została zmieniona}
	\references{\rf{27}}
\end{usecase}


\begin{usecase}{Ustawienie szczegółowych uprawnień}
	% Opis przypadku użycia
	\desc{Twórca Projektu ustawia szczegółowe uprawnienia innym Członkom Projektu}
	% Priorytet
	\priority{6}
	% Aktorzy
	\actors{Użytkownik}
	% Przebieg główny
	\begin{enumerate}
    \item Użytkownik wchodzi w sekcję ‘swoje projekty’
    \item Użytkownik wybiera odpowiedni projekt
    \item Użytkownik wybiera opcję ‘zarządzaj Członkami Projektu’
    \item Użytkownik wybiera odpowiedniego Członka Projektu
    \item Użytkownik zmienia uprawnienia wybranego Członka Projektu
	\end{enumerate}
	% Przebiegi alternatywne
	\beginalternatives
	\alternative{Alternative flow 1}{Działanie się nie wykona, jeżeli użytkownik nie będzie miał uprawnień do nadawania uprawnień innym aktorom.}
	% Diagram
	\diagram{ustawienia_szczegolowych_uprawnien.pdf}{1.0} % drugi argument to szerokość obrazu (1.0 = szerokość strony)
	% Diagram sekwencji
	%\sekwencja{todo.pdf}{0.6}  % drugi argument to szerokość obrazu (1.0 = szerokość strony)
	% Diagram sekwencji
	%\klasa{todo.pdf}{0.6}  % drugi argument to szerokość obrazu (1.0 = szerokość strony)
	% Pre-conditions
	\precond{Użytkownik jest zalogowany i jest Twórcą przynajmniej jednego Projektu }
	% Post-conditions
	\postcond{Uprawnienia co najmniej jednego użytkownika zostały zmienione}
	% Referencje
	\references{\rf{1} \rf{19}}
\end{usecase}

\begin{usecase}{Tworzenie projektu z wieloma zgłoszeniami}
	% Opis przypadku użycia
	\desc{Użytkownik tworzy projekty zawierające wiele zgłoszeń}
	% Priorytet
	\priority{4}
	% Aktorzy
	\actors{Użytkownik}

	% Przebieg główny
	\begin{enumerate}
    \item Użytkownik wybiera opcję ‘utwórz nowy Projekt’
    \item Użytkownik wypełnia odpowiednia informacje o Nowym Projekcie
    \item Użytkownik dodaje innych użytkowników jako Członków Projektu
    \item Użytkownik tworzy Zgłoszenia w danym projekcie.
	\end{enumerate}
	% Przebiegi alternatywne
	\beginalternatives
	\alternative{Alternative flow 1}{Przy odpowiednio dużej ilości zgłoszeń system nie pozwoli tworzenie projektu}
	\alternative{Alternative flow 2}{Jeżeli użytkownik poda niepoprawne hasło system zwróci komunikat o tym}
	% Diagram
	\diagram{tworzenie_projektu_z_wieloma_zgloszeniami.pdf}{1.0} % drugi argument to szerokość obrazu (1.0 = szerokość strony)
	% Diagram sekwencji
	%\sekwencja{todo.pdf}{0.6}  % drugi argument to szerokość obrazu (1.0 = szerokość strony)
	% Diagram sekwencji
	%\klasa{todo.pdf}{0.6}  % drugi argument to szerokość obrazu (1.0 = szerokość strony)
	% Pre-conditions
	\precond{Jest Twórcą przynajmniej jednego Projektu}
	% Post-conditions
	\postcond{Powstanie co najmniej jednego zgłoszenia}
	% Referencje
	\references{\rf{11} \rf{14}}
\end{usecase}

\begin{usecase}{Udzielenie informacji zwrotnej}
	% Opis przypadku użycia
	\desc{Odbiorca Zgłoszenia udziela informacji zwrotnej Osobie Odpowiedzialnej przez zmianę statusu Zgłoszenia.}
	% Priorytet
	\priority{3}
	% Aktorzy
	\actors{Użytkownik}
	% Przebieg główny
	\begin{enumerate}
    \item Odbiorca Zgłoszenia  loguje się
    \item Odbiorca Zgłoszenia wchodzi w sekcję "moje projekty"
    \item Odbiorca Zgłoszenia wybiera Projekt
    \item Odbiorca Zgłoszenia wybiera Zgłoszenie
    \item Odbiorca Zgłoszenia edytuje Zgłoszenie
    \item Odbiorca Zgłoszenia ustawia odpowiedni status Zgłoszenia i zamieszcza komentarz
    \item Odbiorca Zgłoszenia zatwierdza zmiany
	\end{enumerate}
	% Przebiegi alternatywne
	\beginalternatives
	\alternative{Alternative flow 1}{Działanie się nie powiedzie, jeżeli użtkownik będzie miał status projektu zamkniętego}
	% Diagram
	\diagram{udzial_informacji_zwrotnej.pdf}{1.0} % drugi argument to szerokość obrazu (1.0 = szerokość strony)
	% Diagram sekwencji
	%\sekwencja{todo.pdf}{0.6}  % drugi argument to szerokość obrazu (1.0 = szerokość strony)
	% Diagram sekwencji
	%\klasa{todo.pdf}{0.6}  % drugi argument to szerokość obrazu (1.0 = szerokość strony)
	% Pre-conditions
	\precond{Odbiorca Zgłoszenia jest twórcą przynajmniej jednego Projektu}
	% Post-conditions
	\postcond{Zmiana statusu}
	% Referencje
	\references{\rf{11} \rf{13} \rf{18} \rf{20}}
\end{usecase}

\begin{usecase}{Aktualizacja zasobów}
	% Opis przypadku użycia
	\desc{Odbiorca Zgłoszenia aktualizuje Zasoby Zgłoszenia.}
	% Priorytet
	\priority{8}
	% Aktorzy
	\actors{Użytkownik}
	% Przebieg główny
	\begin{enumerate}
    \item Odbiorca Zgłoszenia wchodzi na stronę
    \item Ta osoba loguje się jako root
    \item Próbuje zmienić “propozycję”
    \item Jeżeli taka jest to wraca do punktu 3
    \item Zmiana się powiodła
	\end{enumerate}
	% Przebiegi alternatywne
	\beginalternatives
	\alternative{Alternative flow 1}{Jeżeli po zalogowaniu nie będzie "propozycji" to działanie się nie powiedzie}
	% Diagram
	\diagram{aktualizacja_zasobow.pdf}{1.0} % drugi argument to szerokość obrazu (1.0 = szerokość strony)
	% Diagram sekwencji
	%\sekwencja{todo.pdf}{0.6}  % drugi argument to szerokość obrazu (1.0 = szerokość strony)
	% Diagram sekwencji
	%\klasa{todo.pdf}{0.6}  % drugi argument to szerokość obrazu (1.0 = szerokość strony)
	% Pre-conditions
	\precond{Odbiorca Zgłoszenia jest twórcą przynajmniej jednego Projektu}
	% Post-conditions
	\postcond{Zmiana zasobów zgłoszenia}
	% Referencje
	\references{\rf{17} \rf{18} \rf{19}}
\end{usecase}

\begin{usecase}{Pobieranie załączników}
	% Opis przypadku użycia
	\desc{Twórca/Odbiorca Zgłoszenia/Odbiorca pobiera załączniki dodane do Zgłoszeń.}
	% Priorytet
	\priority{2}
	% Aktorzy
	\actors{Użytkownik, Twórca projektu}
	% Przebieg główny
	\begin{enumerate}
    \item Osoba wchodzi na stronę
    \item Loguje się podając Imię Nazwisko i nr. PESEL
    \item Osoba wyszukuje zgłoszenie
    \item Pobiera załącznik.
    \item Jeżeli to jest Twórca projektu, jest opcja pobrać wszystkie załączniki
	\end{enumerate}
	% Przebiegi alternatywne
	% Diagram
	\diagram{pobieranie_zalacznikow.pdf}{1.0} % drugi argument to szerokość obrazu (1.0 = szerokość strony)
	% Diagram sekwencji
	%\sekwencja{todo.pdf}{0.6}  % drugi argument to szerokość obrazu (1.0 = szerokość strony)
	% Diagram sekwencji
	%\klasa{todo.pdf}{0.6}  % drugi argument to szerokość obrazu (1.0 = szerokość strony)
	% Pre-conditions
	\precond{Odbiorca Zgłoszenia jest twórcą przynajmniej jednego Projektu (\ref{tworzenie_projektu})}
	% Post-conditions
	\postcond{Zmiana zasobów zgłoszenia}
	% Referencje
	\references{\rf{31}}
\end{usecase}

\begin{usecase}{Przeszukiwanie zgłoszeń}
	% Opis przypadku użycia
	\desc{Członek Projektu przeszukuje Zgłoszenia w Projekcie.}
	% Priorytet
	\priority{7}
	% Aktorzy
	\actors{Użytkownik}
	% Przebieg główny
	\begin{enumerate}
    \item Użytkownik loguje się
    \item W polu “Znajdź zgłoszenie” wpisuje kluczowe słowa
    \item Przeglądarka wyświetla znalezione zgłoszenia
    \item Jeżeli nie znalazł nic wyświetla komunikat “Takich zgłoszeń nie ma”
	\end{enumerate}
	% Przebiegi alternatywne
	\beginalternatives
	\alternative{Alternative flow 1}{Jeżeli nie będzie znaleziono żadnego ze zgłoszeń aplikacja wyświetli odpowiedni komunikat}
	% Diagram
	\diagram{przeszukiwanie_zgloszen.pdf}{1.0} % drugi argument to szerokość obrazu (1.0 = szerokość strony)
	%\sekwencja{todo.pdf}{0.6}  % drugi argument to szerokość obrazu (1.0 = szerokość strony)
	% Diagram sekwencji
	%\klasa{todo.pdf}{0.6}  % drugi argument to szerokość obrazu (1.0 = szerokość strony)
	% Pre-conditions
	\precond{Odbiorca Zgłoszenia jest twórcą przynajmniej jednego Projektu (\ref{tworzenie_projektu})}
	% Post-conditions
	\postcond{Uzyskanie odpowiednich zgłoszeń lub informacji że ich brak}
	% Referencje
	\references{\rf{32}}
\end{usecase}

\begin{usecase}{Wysyłanie informacji o zamknięciu projektu, "Odbiorca-Twórca"}
	% Opis przypadku użycia
	\desc{Odbiorca Zgłoszenia wysyłam Twórcy Zgłoszenia informacje, że zgłoszenie jest do zamknięcia}
	% Priorytet
	\priority{5}
	% Aktorzy
	\actors{Użytkownik, Twórca projektu}
	% Przebieg główny
	\begin{enumerate}
    \item Odbiorca Zgłoszenia loguje się do systemu
    \item Odbiorca Zgłoszenia wchodzi do sekcji zgłoszenia
    \item W tym zgłoszenie jest podany adres @mail
    \item Odbiorca Zgłoszenia wysyła mail
	\end{enumerate}
	% Przebiegi alternatywne
	\beginalternatives
	\alternative{Alternative flow 1}{Jeżeli adres @mail nie będzie podany,
	system wyświetli odpowiedni komunikat}
	% Diagram
	\diagram{odbiorca_tworca.pdf}{1.0} % drugi argument to szerokość obrazu (1.0 = szerokość strony)
	%\sekwencja{todo.pdf}{0.6}  % drugi argument to szerokość obrazu (1.0 = szerokość strony)
	% Diagram sekwencji
	%\klasa{todo.pdf}{0.6}  % drugi argument to szerokość obrazu (1.0 = szerokość strony)
	% Pre-conditions
	\precond{Odbiorca i twórca muszą należeć do co najmniej jednego projektu.}
	% Post-conditions
	\postcond{Uzyskanie odpowiednich zgłoszeń lub informacji, że jest ich brak}
	% Referencje
	\references{\rf{20} \rf{21}}
\end{usecase}

\begin{usecase}{Wysłanie informacji o zamknięciu projektu, "Użytkownik-Twórca"}
	% Opis przypadku użycia
	\desc{Członek Projektu wysyła Twórcy Projektu informacje, że  projekt jest do zamknięcia.}
	% Priorytet
	\priority{5}
	% Aktorzy
	\actors{Użytkownik, Twórca projektu}

	% Przebieg główny
	\begin{enumerate}
    \item Członek projektu wchodzi na stronę
    \item Podaje Imię, Nazwisko i identyfikator
    \item Naciska na przycisk “Projekt jest do zamknięcia”
    \item Wiadomość przychodzi do Twórcy projektu
	\end{enumerate}
	% Przebiegi alternatywne
	\beginalternatives
	\alternative{Alternative flow 1}{Jeżeli adres @mail nie będzie podany,
		system wyświetli odpowiedni komunikat}
	% Diagram
	\diagram{uzytkownik_tworca.pdf}{1.0} % drugi argument to szerokość obrazu (1.0 = szerokość strony)
	%\sekwencja{todo.pdf}{0.6}  % drugi argument to szerokość obrazu (1.0 = szerokość strony)
	% Diagram sekwencji
	%\klasa{todo.pdf}{0.6}  % drugi argument to szerokość obrazu (1.0 = szerokość strony)
	% Pre-conditions
	\precond{Odbiorca i twórca muszą należeć do co najmniej jednego projektu.}
	% Post-conditions
	\postcond{Uzyskanie odpowiednich zgłoszeń lub informacji że jest ich brak}
	\references{\rf{20} \rf{21}}
\end{usecase}

\chapter{Wymagania funkcjonalne}
W niniejszym rozdziale przedstawione zostały wymagania funkcjonalne, które zdefiniowano dla poszczególnych modułów systemu. Priorytet oraz złożoność realizacji każdego z wymagań wyrażone zostały przy użyciu dziesięciostopniowej skali liniowej (1--10). Liczba 10 odpowiada najwyższemu priorytetowi i największej złożoności, natomiast liczba 1 - wartościom najniższym.
\section{Moduł uwierzytelniania}
% Moduł uwierzytelniania

\funcreq{System powinien pozwolić gościom zarejestrować się.}{10}{4}
\funcreq{System powinien pozwolić użytkownikom zalogować i wylogować się.}{10}{5}
\funcreq{System powinien pozwolić użytkownikom zmieniać swoje dane (takie jak adres e-mail, hasło).}{7}{4}
\funcreq{System powinien weryfikować dane wprowadzane przez użytkowników (takie	jak typ danych lub czy użytkownik o danej nazwie / adresie e-mail już nie istnieje).}{10}{5}
\section{Moduł administracyjny}
% Moduł administracyjny

\funcreq{System powinien umożliwiać Administratorom weryfikację nowo-zarejestrowanych użytkowników.}{10}{6}
\funcreq{System powinien umożliwiać Administratorom zmianę uprawnień innych użytkowników.}{10}{4}
\funcreq{System powinien umożliwiać Administratorom usuwanie kont innych użytkowników.}{9}{4}
\funcreq{System powinien umożliwiać Administratorom robienie backupu danych.}{10}{7}
\funcreq{System powinien umożliwiać Administratorom przeglądanie logów systemu.}{10}{3}
\funcreq{System powinien umożliwiać Administratorom przeprowadzenie rekonfiguracji Systemu.}{10}{8}
\section{Moduł główny}
% Moduł główny

\funcreq{System powinien pozwolić użytkownikom na tworzenie nowego projektu. \label{3_1}}{10}{5}
\funcreq{System powinien pozwolić Twórcom Projektu dodawanie Członków Projektu.}{10}{3}
\funcreq{System powinien pozwolić Członkom Projektu na przeglądanie Projektów do których należą.}{10}{2}
\funcreq{System powinien pozwolić Członkom projektu na tworzenie nowych zgłoszeń.}{10}{5}
\funcreq{System powinien pozwolić Członkom projektu przeglądać Zgłoszenia.}{10}{2}
\funcreq{System powinien pozwolić na przypisywanie osób odpowiedzialnych za zgłoszenie pośród Członków projektu}{9}{2}
\funcreq{System powinien pozwolić Osobom odpowiedzialnym za zgłoszenie zmianę jego statusu.}{10}{2}
\funcreq{System powinien pozwolić Twórcom zgłoszenie weryfikację jego realizacji i odpowiednią zmianę statusu.}{9}{2}
\funcreq{System powinien umożliwiać Twórcy projektu edycję swojego projektu.}{10}{4}
\funcreq{System powinien umożliwiać Twórcy projektu zgłoszenie swojego projektu do zamknięcia.}{10}{2}
\funcreq{System powinien umożliwiać Twórcy Zgłoszenia zamknięcie swojego Zgłoszenia.}{10}{2}
\funcreq{System powinien umożliwiać Twórcy Zgłoszenia zmianę widoczności swojego Zgłoszenia.}{5}{2}
\funcreq{System powinien umożliwiać Twórcy Zgłoszenia edycję podstawowych informacji o swoim Zgłoszeniu.}{9}{4}
\funcreq{System powinien umożliwiać Twórcy Zgłoszenia wyznaczenie Użytkownika będącego Osobą Odpowiedzialną za Zgłoszenie.}{9}{3}
\funcreq{System powinien umożliwiać Odbiorcy Zgłoszenia udzielenie informacji zwrotnej Osobie Odpowiedzialnej za Zgłoszenie.}{9}{3}
\funcreq{System powinien umożliwiać Twórcy Zgłoszenia planowanie kolejnych etapów Projektu w postaci grafu.}{10}{7}
\funcreq{System powinien umożliwiać Członkowi Projektu edycję Grafu poprzez dodawanie, edytowanie i usuwanie Zgłoszeń.}{10}{7}
\funcreq{System powinien umożliwiać Osobie Odpowiedzialnej aktualizowanie Zasobów Zgłoszenia.}{8}{5}
\funcreq{System powinien umożliwiać Osobie Odpowiedzialnej wysyłanie do Twórcy Zgłoszenia informacji, że Zgłoszenie jest do zamknięcia.}{6}{2}
\funcreq{System powinien umożliwiać Członkowi Projektu wysyłanie do Twórcy Projektu informacji, że Projekt jest do zamknięcia.}{6}{2}
\funcreq{System powinien umożliwiać Twórcy, Osobie Odpowiedzialnej, Odbiorcy pobieranie załączników dodanych do Zgłoszeń.}{4}{5}
\funcreq{System powinien umożliwiać Członkowi Projektu przeszukiwanie Zgłoszeń w Projekcie.}{8}{5}

\chapter{Wymagania niefunkcjonalne}
\section{Dostępność}
\nonfuncreq{System powinien być dostępny minimum 99\% czasu.}{10}{8}
\nonfuncreq{System powinien być skonstruowany w sposób umożliwiający jego szybką naprawę}{\todo}{\todo}
\section{Bezpieczeństwo}
\nonfuncreq{System powinien robić backup bazy danych codziennie podczas małego obciążenia (prawdopodobnie o 3:15 w nocy). Nie powinno to powodować przerwy w dostępności.}{10}{6}
\nonfuncreq{Połączenie z serwerem jest szyfrowane.}{10}{6}
\nonfuncreq{System powinien generować logi zawierające informacje o zmianach dokonywanych przez użytkowników oraz informacje diagnostyczne.}{7}{4}
\nonfuncreq{System powinien wykrywać próby nieautoryzowanego dostępu i informować o nich administratora.}{7}{5}
\nonfuncreq{System powinien być blokowany na 10 min dla IP adresu z którego próbowano się połączyć więcej niż 5 razy.}{8}{6}
\section{Skalowalność}
\nonfuncreq{System powinien być skalowalny (dwukrotne zwiększenie ilości serwerów powinno pozwolić na dwukrotne zwiększenie ilości klientów/ kont użytkowników).}{7}{5}
\nonfuncreq{System powinien być elastyczny, czyli ma mieć możliwość zmieniać swoje właściwości ze zmianą popytu na produkcję.}
\section{Ergonomia}
\nonfuncreq{System powinien być ergonomiczny, czyli wygodny w użyciu.}{10}{9}
\nonfuncreq{Wygląd Systemu powinien być estetyczny, czyli wygodny w użyciu oraz zachęcający nowych Użytkowników.}{10}{9}


\chapter{Diagramy sekwencji}
\section{Moduł uwierzytelniania}

\begin{sequence}{Logowanie}
\diagramSQ{logowanie_sekwencja.pdf}{1}
\end{sequence}

\begin{sequence}{Rejestracja}
\diagramSQ{rejestracja_sekwencja.pdf}{1}
\end{sequence}

\section{Moduł administracyjny}

\begin{sequence}{Zatwierdzanie konta}
\diagramSQ{zatwierdzenie_konta_sekwencja.pdf}{1}
\end{sequence}

\begin{sequence}{Usuwanie konta}
\diagramSQ{usuwanie_konta_sekwencja.pdf}{1}
\end{sequence}

\begin{sequence}{Zmiana uprawnień}
\diagramSQ{zmiana_uprawnien_sekwencja.pdf}{1}
\end{sequence}

\begin{sequence}{Backup danych}
\diagramSQ{backup_sekwencja.pdf}{1}
\end{sequence}

\begin{sequence}{Przywracanie systemu}
\diagramSQ{przywracanie_systemu_sekwencja.pdf}{1}
\end{sequence}

\begin{sequence}{Przeglądanie logów Systemu}
\diagramSQ{przegladanie_logow_sekwencja.pdf}{1}
\end{sequence}

\begin{sequence}{Rekonfiguracja Systemu}
\diagramSQ{rekonfiguracja_sekwencja.pdf}{1}
\end{sequence}

\section{Moduł główny}

\begin{sequence}{Przeglądanie Projektów}
	\label{Przegladanie projektow}
	\diagramSQ{przegladanie_projektow_sekwencja_fit.pdf}{0.9}
\end{sequence}

\begin{sequence}{Edycja Projektu}
	\label{Edycja projektu}
	\diagramSQ{edycja_projektu_sekwencja_fit.pdf}{0.9}
\end{sequence}

\begin{sequence}{Zamknięcie projektu}
	\label{Zamkniecie projektu}
	\diagramSQ{zamkniecie_projektu_sekwencja_fit.pdf}{0.9}
\end{sequence}

\begin{sequence}{Tworzenie zgłoszenia}
	\label{Tworzenie zgloszenia}
	\diagramSQ{tworzenie_zgloszenia_sekwencja_fit.pdf}{0.9}
\end{sequence}

\begin{sequence}{Zamknięcie zgłoszenia}
	\label{Zamkniecie zgloszenia}
	\diagramSQ{zamkniecie_zgloszenia_sekwencja_fit.pdf}{0.9}
\end{sequence}

\begin{sequence}{Zmaina widocznosci zgłoszenia}
	\label{Zmiana widocznosci zgloszenia}
	\diagramSQ{zmiana_widocznosci_zgloszenia_sekwencja_fit.pdf}{0.9}
\end{sequence}

\begin{sequence}{Edycja zgłoszenia}
	\label{Edycja zgloszenia}
	\diagramSQ{edycja_zgloszenia_sekwencja_fit.pdf}{0.9}
\end{sequence}




\chapter{Diagramy Klas}
\section{Moduł uwierzytelniania}

\begin{class}{Logowanie}
\diagramSQ{logowanie_klas.pdf}{1}
\end{class}

\begin{class}{Rejestracja}
	\diagramSQ{rejestracja_klas.pdf}{1}
\end{class}

\section{Moduł administracyjny}

\begin{class}{Zatwierdzanie konta}
	\diagramSQ{zatwierdzenie_konta_klas.pdf}{1}
\end{class}

\begin{class}{Usuwanie konta}
	\diagramSQ{usuwanie_konta_zmiana_uprawnien_klas.pdf}{1}
\end{class}

\begin{class}{Backup}
	\diagramSQ{backup_klas.pdf}{1}
\end{class}

\begin{class}{Przywracanie Systemu}
	\diagramSQ{przywracanie_systemu_klas.pdf}{1}
\end{class}

\begin{class}{Przeglądanie Logów}
	\diagramSQ{przegladanie_logow_klas.pdf}{1}
\end{class}

\begin{class}{Zmiana uprawnień}
	\diagramSQ{usuwanie_konta_zmiana_uprawnien_klas.pdf}{1}
\end{class}

\begin{class}{Rekonfiguracja Systemu}
	\diagramSQ{rekonfiguracja_klas.pdf}{1}
\end{class}

\section{Moduł główny}



\chapter{Diagram komponentów}
%jednak bez tego bo bedzie ogolny jeden, wiec nie trzeba liczyc
%\begin{component}{Diagram komponentów całego systemu}
\diagramC{diagram_komponentow.png}{1}
%\end{component}




\end{document}
